\section{Why the result is full iPRA bi-interpretability}
\label{sec:bridge}

It would be insufficient to show that the new rules hold over the standard
natural numbers. It would also be insufficient to recover the equations of
primitive recursive functions while ignoring quantified intuitionistic
formulas. Here the new comparison preserves the entire open consequence
relation, and the remaining interpretation chain is the one already
established in the pinned sources.

\subsection{The guarded/native comparison}

Let $\mathcal C$ be the unrestricted Choice-DI calculus: it has raw scalar
terms, Eq/$D$ formulas, the two computation schemas, Use, full formula Tr,
Open, case introduction at every matrix, and progress introduction at literal
equations. Its ordinary reading is the one stated for $D$ above.

For an equation $E=(s=t)$ define
\[
 G_vE=(I(v,0;u.s)=I(v,0;u.t)),\qquad
 L(E)=D_{a,v}G_vE,
\]
where the added variables are fresh. Define $P(A)=L(A)$ when $A$ is an
equation and $P(A)=A$ otherwise. Then $P(A)$ is always compound. The
guarded-core proof~\cite{guarded} establishes, under arbitrary retained
contexts,
\[
 L(E)\dashv\vdash E.
\]
Its forward direction opens the clause, supplies the true zero row, and
specializes at $1$. Its reverse direction uses the retained $E$ to prove the
successor row, ignoring the induction hypothesis.

The structural translation into the guarded language is
\[
 T(s=t)=(s=t),\qquad
 T(D_{x,y}B)=D_{x,y}P(TB).
\]
Literal inclusion is the interpretation in the other direction. Case
introductions translate through padding. Equation progress translates using
the guarded calculus's derived equation-induction rule. An opened padded
clause is used by padding its antecedent and unpadding its consequence.
These are transformations of arbitrary finite open derivations.

The two formula composites are also proved, not assumed. In the unrestricted
calculus, padding and general binder congruence give $A\dashv\vdash JT(A)$.
In the guarded calculus, compound matrices use compound binder congruence;
guarded atomic matrices use their provable zero row to establish
\[
 D_{x,y}B\dashv\vdash D_{x,y}L(B).
\]
This is why the zero guard in the public grammar is substantive. Dropping it
would invalidate this argument. The source proof, supplied in
\path{reference/guarded_core.tex}, gives the equation-induction bootstrap and
both composites in detail. The new deterministic calculus proves exactly the
same guarded sequents, so every one of those derivations transfers literally.

\subsection{The ordinary scalar theory}

The intermediary $T_{\rm eq}$ uses ordinary intuitionistic first-order logic
with equality over scalar terms. Its nonlogical principles are
\[
 I(0,a;u.h)=a,\quad I(Sn,a;u.h)=h[I(n,a;u.h)/u],
\]
\[
 \forall x\,(x=0\lor\exists y\,x=Sy),\qquad 0=S0\to\bot,
\]
together with contextual induction for literal equations. ``Contextual''
means the retained side assumptions may themselves be quantified formulas;
only the induction formula is restricted to an equation.

The ordinary realization sends $D_{x,y}B$ to
$\exists x(B(x,0)\to\forall y\,B(x,y))$ and sends a clause to
$A\to\forall y B$. In $T_{\rm eq}$, compound case introduction is justified
by coverage; guarded atomic progress is justified by equation induction and
the true zero row. Equality transport and clause application use the usual
intuitionistic equality and quantifier rules. Thus the reading is sound.

Conversely, the native logical macros recover implication, conjunction,
disjunction, both quantifiers, and falsity, with their introduction and
elimination rules, and prove the formula round trips. This is the
full native/ordinary bridge~\cite{choice,ordinary,lean}. No excluded-middle
rule for arbitrary formulas is used. The code of a primitive recursive
description is never an object-language numerical variable ranging over
all programs.

\subsection{The arithmetic interpretation and its inverses}

In iPRA, each fixed scalar term has a primitive recursive interpretation.
An occurrence $I(n,a;u.h)$ becomes the primitive recursion that starts at the
translation of $a$ and repeatedly applies the already translated update $h$,
ignoring the recursion counter in that update. This is external recursion
over a fixed finite term description, not an internal universal evaluator.

In the reverse direction the repository develops, in order, a scalar
predecessor, reflection of actual equality into a Boolean-valued term,
quantifier-free induction, pairing with projections, and a compiler for
arbitrary fixed primitive-recursive descriptions~\cite{arithmetic}.
For a recursion
\[
 f(0,\vec p)=g(\vec p),\qquad
 f(Sn,\vec p)=h(n,f(n,\vec p),\vec p),
\]
the scalar compiler stores the counter and current output in one coded state.
Writing $\pi$ for pairing and $\ell,r$ for its two projections, put
\[
 W(u,\vec p)=\pi(S\ell(u),h^*(\ell(u),r(u),\vec p)),
\]
\[
 V(n,\vec p)=I(n,\pi(0,g^*(\vec p));u.W(u,\vec p)),
 \qquad f^*(n,\vec p)=r(V(n,\vec p)).
\]
The invariant $\ell(V(n,\vec p))=n$, proved by equation induction,
and the projection equations prove the two defining equations internally.
Initial functions and composition translate directly.

Equality reflection compiles each quantifier-free formula $A$ to a fixed
term $b_A$ with $b_A=0\leftrightarrow A$ and a proved Boolean range. Equation
induction on $b_A=0$ therefore yields full quantifier-free induction,
including under arbitrary side assumptions. This is an inherited bootstrap;
we do not use equality reflection to justify a restricted induction rule
before reflection has been proved.

Both numerical interpretation domains are $x=x$. The interpreted equality
is actual equality, and the numerical isomorphisms for the composites have
graph $x=y$. Internal uniqueness of a fixed primitive recursion, proved by
equation induction, identifies each double-translated operation with the
original one. Structural induction on each fixed description then proves
preservation of every operation; the native formula round trips give the
logical composites as well. These claims concern arbitrary internal numbers,
not merely standard numerals.

The precise downstream Lean statements in the pinned source include
\path{native_iff_pra}, \path{pra_iff_native},
\path{native_roundTrip}, \path{pra_native_roundTrip}, and
\path{native_inconsistent_iff} in
\path{Lean/ConciseIPRA/NativeBiInterpretation.lean}. This package includes
that source for reference. It does not claim that this file alone is a
standalone Lean project, or that the new deterministic-introduction theorem
has been compiled in Lean.

\section{A failed shortcut and the boundary it exposes}

The current automatic-opening proposals inspect a zero row and a first row.
One tempting simplification is to keep only the first-row check. It is
unsound. This matters when trying to combine the new introduction rule with
those proposals.

Choose a primitive recursive $f$ with
\[
 f(0)=f(1)=0,\qquad f(2)=1,
\]
and let $B(x,y)$ be $x=f(y)$. In full iPRA,
\[
 D_{x,y}B=\exists x\bigl(x=f(0)\to\forall y\,x=f(y)\bigr)
\]
is provable with witness $1$: its antecedent is false. The first row at a
fresh parameter $b$ is $b=0$. An unchecked first-row probe would therefore
replace the opened witness by $0$. The resulting clause is
\[
 0=0\to\forall y\,0=f(y),
\]
which gives $0=1$ at $y=2$.

Our literal naming matrix has a computation-provable zero row and is used
with ordinary Open; it does not make this invalid inference. Nor do we
silently compose it with automatic opening: opening a successful name
directly at $t$ can remove the fresh alias $b=t$ on which the present
simulation relies. A composition needs an additional proof, rather than a
simple replacement of one rule name by another.

\section{Further research questions}

\begin{enumerate}
\item \textbf{Can automatic introduction and automatic opening share a
      smaller bootstrap?} The present proof uses full formula transport and
      ordinary opening. A useful next target is a simultaneous recovery of
      transport, naming, and deterministic introduction in the three-family
      full-normal-form calculus. The dependence of existing equality proofs
      on unrestricted introduction must be removed explicitly.
\item \textbf{Can the deterministic selector be made cheaper to recognize?}
      Its mathematical data are smaller, but a pattern recognizer can exceed
      the source cost of validating a supplied term. A redesigned formula
      grammar may make the name structure intrinsic; any extra grammar and
      translation costs must be counted.
\item \textbf{Can full formula Tr be weakened before the bootstrap?}
      The algebraic-cut argument uses transport through an arbitrary guarded
      formula. Reducing it to term-only transport would require proving
      formula transport without already using unrestricted witnesses or
      unrepaired weakening.
\item \textbf{What is the proof-size cost?} The construction is effective,
      but this report does not claim a linear or polynomial bound on the
      serialized certificates. An explicit compiler could measure duplication
      caused by repaired weakening and naming.
\item \textbf{Can the new structural theorem be formalized against the
      existing Lean kernel?} A direct port using the substitution-stable
      diagonal presentation is plausible, but the guarded presentation and
      its exact source interpretation should first be pinned in the formal
      interface. A source-level prototype is not a compiled theorem.
\item \textbf{Which size metric captures the intended simplicity?}
      Rule-data count, complete typed-checker size, compact encoding size,
      and conceptual prerequisites reward different changes. A future
      minimum claim should fix one public input contract and one measured
      implementation language, and distinguish observed records from lower
      bounds.
\end{enumerate}

